% Suggested replacement for the converse paragraph in the proof of Lemma
% \ref{lem-kraft} in chapters-tex/shannon-coding.tex. Uses the book's notation.
% This fragment is not included in the published book.

$\Leftarrow$ Conversely, suppose that $\sum_k2^{-\ell_k}\leq1$.
Relabel the symbols so that $\ell_1\leq\cdots\leq\ell_K=:m$.
In the complete binary tree of depth $m$, number the leaves from left to
right by $0,\ldots,2^m-1$, according to their values as binary integers.
A word of length $\ell$ representing the integer $q$ has the descendant
leaves
\[
 \{q2^{m-\ell},\ldots,(q+1)2^{m-\ell}-1\}:
\]
appending $m-\ell$ bits amounts to adding an integer between $0$ and
$2^{m-\ell}-1$ to $q2^{m-\ell}$. Thus a consecutive block of
$2^{m-\ell}$ leaves is the leaf set of a subtree rooted at depth $\ell$
precisely when its starting index is divisible by $2^{m-\ell}$.

Set $b_k=2^{m-\ell_k}$ and $S_k=\sum_{j=1}^k b_j$, with $S_0=0$.
We assign the $k$th symbol the block
$B_k=\{S_{k-1},\ldots,S_k-1\}$. These blocks are disjoint and fit among
the $2^m$ leaves because
\[
 S_K=2^m\sum_k2^{-\ell_k}\leq2^m.
\]
Moreover, their starting indices have exactly the required alignment:
\[
 q_k:=\frac{S_{k-1}}{b_k}
     =\sum_{j<k}2^{\ell_k-\ell_j}\in\ZZ_{\geq0},
\]
since $\ell_j\leq\ell_k$ for $j<k$. The inequality
$(q_k+1)b_k=S_k\leq2^m$ gives $q_k<2^{\ell_k}$. Write $q_k$ in binary
using exactly $\ell_k$ bits, including leading zeroes, and call the resulting
word $c_k$. Its depth-$m$ descendant leaves are precisely $B_k$.

If $c_j$ were a prefix of $c_k$ for $j\ne k$, the nonempty descendant block
$B_k$ would be contained in $B_j$, contradicting disjointness. Hence the
words are prefix-free and have the prescribed lengths. The union of their
root-to-codeword paths is a binary tree whose leaves are exactly the
$c_k$. Undo the initial relabeling to recover the original symbol order.
If some $\ell_k=0$, the inequality forces $K=1$, and the code consists of
the empty word.
